\subsection{Native authentication for the selected recovered root}
\label{sec:native-authentication}

The confirmation and first-message authentication bound can be transported
through the native execution by comparing its own terminal event bit.  This
gives an end-to-end authentication statement for the same recovered root whose
first message appears in the confidentiality game.

\paragraph{Target event.}
Fix the live-core, single-root next-rekey profile and an admitted marked
controller of Theorem~\ref{thm:cb-next-rekey} and
Proposition~\ref{prop:confirmation-message-auth}.  In particular, for each
fixed candidate its target contribution and marked traffic entry are
designated at the allocation points required by those interfaces; the
original terminal selector corresponds to those designated entries whenever
it returns that eligible candidate.  As in
Definition~\ref{def:recovery-game} and its adaptive version, the experiment
retains one application challenge.  Its original accepted logical send fixes
the target root schedule through the send's stored provenance handle.  The
terminal selector may choose coordinates, reservation handles, and the cut
adaptively; eligibility requires those coordinates to identify that same
accepted send.  Let $E$ be its original eligibility event, $A$ its
source-attribution event, and $I\in J=[G]\times[C]$ the attributed pair of
original generation and encapsulation reservations on $A$.

The target root and its event history are defined by the actual accepted
challenge send, including on executions where $A$ fails.  A challenge without
an accepted send has $E=0$ and we set the target event flag to zero.  On $A$,
the root handle is the schedule entry generated from the effective
ContextBound tuple $m^*$ of the attributed pair.  The handle is carried
through replay, recovery, exact tuple aliases, and the resulting role keys
and marked first-message key.  Ledger aliasing follows
Proposition~\ref{prop:confirmation-message-auth}.

For this target, let $\mathsf{Bad}_{\rm auth}$ be the union of the following
events: a successful native response, final, or receipt confirmation whose
complete core is absent from the earlier honest private-tag ledger under the
corresponding target role key; or a fresh authenticated first-message receive
and inbox installation whose $(\mathrm{AAD},\mathrm{ciphertext})$ pair is absent
from the earlier honest private-encryption ledger under the marked target
first-message key.  These are the events of the uniform-root authentication
proposition restricted to this selected root schedule.  The honest ledgers
are populated at the first private computation, before persistence or public
release.  A repeated delivery or a previously consumed inbox result retains
its original ledger provenance.  Genuine aliases of the selected schedule
entry share its honest ledger, while independently allocated entries retain
their own provenance.

The experiment records these event flags as execution proceeds.  Its eventual
evaluation reads the accepted challenge send's root history, which may
include confirmation events before that send.  For a fixed pair $j=(g,a)$,
the original $a$ determines $m^*$ before its first private hash computation
at any protected alias, including a computation followed by a failed write,
as specified in Section~\ref{sec:contextbound-connection}.  Other protected
tuples retain their own entries.  If the original $a$ never supplies a usable
target, the candidate event is zero.  The root entry, traffic entry, and
direction retain their prescribed online marks.  These rules connect the
recorded history to the admitted controller throughout the comparisons.
The probability bounded below is the joint probability
$\Pr[E\wedge\mathsf{Bad}_{\rm auth}]$ in the original native game.  The
attribution residual bounds its part on $E\wedge\neg A$.

\paragraph{Authentication-specific comparisons.}
For each fixed $j\in J$, use the native stopped continuation
$G_{\rm nat}^{[j]}$, its real-component record idealization $V^{[j]}$, and the
uniform-component continuation $\widetilde V^{[j]}$ from the candidate-local
accounting.  Apply to each of these worlds the same terminal Boolean code
\[
 W_j=\mathbf1[
 E\wedge A\wedge I=j\wedge\mathsf{Bad}_{\rm auth}],
 \qquad W_j=0\text{ after a declared local stop}.
\]
All eligibility, attribution, tuple selection, ledger membership, and event
flags are evaluated on the execution of that world.  They are computable from
the classical execution and provenance ledgers available to the existing
reduction interface; no test of an unknown challenge bit is used.

Let $a_j^{\rm auth}$ be the real-versus-uniform advantage of the operational
ML-KEM adapter of Section~\ref{sec:contextbound-connection} when its output
bit is $W_j$.  Its public key is installed at the original generation's
first actual private invocation, its encapsulation challenge at the original
encapsulation's first actual private invocation, and every post-challenge decapsulation of
the selected PQ ciphertext under the selected generation owner returns the
challenge secret, regardless of the traditional ciphertext or context.
The native traditional-component validation, its rejections, the original
prechallenge history, and exact-retry caches are retained.  The selected-pair correctness allowance is
$\delta_j$.  Consequently this adapter gives
\begin{equation}
 \left|\Pr_{V^{[j]}}[W_j=1]
       -\Pr_{\widetilde V^{[j]}}[W_j=1]\right|
 \le a_j^{\rm auth}+\delta_j.
 \label{eq:native-auth-component}
\end{equation}
This uses the same two-stage component experiment as the existing adapter,
with the authentication event as its output bit.

Use the whole-continuation budgets of the ContextBound, root-initialization,
joint-schedule, and confirmation comparisons.  Define
\begin{align}
 \Omega_j^{\rm auth}
 :=\min\bigl\{1,{}&
 \Phi_{256}(d_j,Q_{H,j})+\Delta_{{\rm init},j}
       +\Delta_{{\rm sched},j}+\Delta_{{\rm conf},j}
       +V_j2^{-256}\nonumber\\[-1mm]
       &+\chi_H(Q_{H,j}^{\rm conf})
       +\varepsilon_{{\rm INT},j}\bigr\}.
 \label{eq:native-auth-continuation}
\end{align}
Here $d_j$ counts distinct effective protected tuples in the ContextBound
comparison; $\Delta_{{\rm init},j}$ is the event-bit cost of replacing the
selected contribution's single native root-initialization entry by a fresh
root; $\Delta_{{\rm sched},j}$ is the shared-compression joint-schedule cost;
$\Delta_{{\rm conf},j}$ replaces the three confirmation HMAC rows;
$V_j$ counts the native confirmation verification attempts;
$\chi_H$ is the framed-input collision bound; and
$\varepsilon_{{\rm INT},j}$ is the one-key, at-most-one-encryption
first-message AEAD integrity advantage with the native decryption interface.
Each query budget includes preparation, ordinary primitive computations,
honest protocol work, later continuation, and terminal event evaluation of
its derived experiment.  Budgets are fixed before interaction and are not
reset at a challenge or at a hybrid boundary.

\begin{theorem}[Native selected-root authentication]
\label{thm:native-selected-authentication}
In the live-core single-root next-rekey profile, with the admitted marked
controller just specified and the origin-preservation interfaces of
Section~\ref{sec:native-correspondence}, Lemma~\ref{lem:cb-native-root-use},
and Theorem~\ref{thm:joint-message-leaf},
\begin{equation}
 \Pr_{G_{\rm nat}}[E\wedge\mathsf{Bad}_{\rm auth}]
 \le
 \min\!\left\{1,\;
 \eta_{\rm attr}^{\rm nat}
 +\sum_{j\in J}
  \min\{1,\varepsilon_{{\rm rec},j}
           +a_j^{\rm auth}+\delta_j+\Omega_j^{\rm auth}\}
 \right\}.
 \label{eq:native-selected-authentication}
\end{equation}
The record cost $\varepsilon_{{\rm rec},j}$ is the live-core record event
cost~\eqref{eq:live-core-record-cost}, including its original fresh-seal
equality allowance.  The attribution
residual is measured once in the original native game.  For the executed-slot
selector, it may be instantiated by
Proposition~\ref{prop:executed-source-attribution} and the established
reservation-signature bridge.  An explicitly modeled projection extension adds
its original selected-candidate residual $\eta_{\rm proj}^{\rm nat}$ and uses
the corresponding event-bit record/report comparison in place of
$\varepsilon_{{\rm rec},j}$.
\end{theorem}

\begin{proof}
First localize the native event.  Write
$F=E\wedge\mathsf{Bad}_{\rm auth}$.  On attributed native executions the
events $I=j$ partition $F\wedge A$.  For a general declared projection
extension, split off $F\wedge A\wedge B_I$.  Its probability is at most
$\eta_{\rm proj}^{\rm nat}=\Pr[E\wedge A\wedge B_I]$; likewise
$\Pr[F\wedge\neg A]\le\eta_{\rm attr}^{\rm nat}$.
For each fixed $j$, couple the original execution and $G_{\rm nat}^{[j]}$
until the first free or charged local stop.  The free-stop certificate
permanently falsifies $E\wedge A\wedge I=j$ on every continuation.  A charged
stop is absent on the summand restricted to $\neg B_j$.  Therefore the same
identical-prefix argument as Lemma~\ref{lem:native-local-score}, now applied to a
Boolean variable, gives the exact identity
\[
 \Pr_{G_{\rm nat}}[F\wedge A\wedge I=j\wedge\neg B_j]
 =\Pr_{G_{\rm nat}^{[j]}}[W_j=1].
\]
This identity remains valid when a recorded authentication event occurs
before the challenge send or terminal selection: the stopped predicate is the conjunction
with final eligibility and attribution, and its free-stop certificate
concerns every possible continuation.  Consequently
\begin{equation}
 \Pr_{G_{\rm nat}}[F]
 \le\eta_{\rm attr}^{\rm nat}+\eta_{\rm proj}^{\rm nat}
    +\sum_j\Pr_{G_{\rm nat}^{[j]}}[W_j=1].
 \label{eq:native-auth-localization}
\end{equation}
The live-core interface has $B_j=\varnothing$ and
$\eta_{\rm proj}^{\rm nat}=0$.

Next compare a fixed continuation.  The native representation invariant and
immutable-history simulator preserve the selected event's classical
execution, role-key provenance, private-tag and private-encryption ledgers,
and original allocation indices.  The real-to-opaque record comparison
therefore applies with output $W_j$.  Because it is a bit, its probability
changes by at most $\varepsilon_{{\rm rec},j}$.  This gives
\[
 \Pr_{G_{\rm nat}^{[j]}}[W_j=1]
 \le\varepsilon_{{\rm rec},j}
       +\Pr_{V^{[j]}}[W_j=1]
 \le\varepsilon_{{\rm rec},j}+a_j^{\rm auth}+\delta_j
       +\Pr_{\widetilde V^{[j]}}[W_j=1],
\]
where the second inequality is the operational component comparison
\eqref{eq:native-auth-component}.  All role-key and honest-origin ledgers
remain private reduction bookkeeping.  Membership is checked by recording
the actual complete core or encryption pair at its native invocation; it
requires no inspection of quantum queries or unknown component tape.

In $\widetilde V^{[j]}$, the protected PQ component is an independent uniform
256-bit value.  Apply Lemma~\ref{lem:cb-hidden-component} to the same Boolean
continuation, with the complete effective tuple cache.  It replaces each
protected ContextBound answer by a private random-function entry at cost
$\Phi_{256}(d_j,Q_{H,j})$.  Exact aliases, including aliases occurring before
the eventual selected send, share their entry.  The target entry is marked
at its first actual private use as prescribed by the existing controller
interface, so its fresh value is sampled without conditioning on a future
selection.  If this target is never allocated or never becomes eligible,
$W_j=0$.

Lemma~\ref{lem:cb-native-root-use} places every use of this target contribution
in one old-root and complete-info initialization entry.  Its event-bit
comparison~\eqref{eq:appendix-root-initialization} costs
$\Delta_{{\rm init},j}$ and leaves the uniform selected root controller of
Proposition~\ref{prop:confirmation-message-auth}.  All other tuple
entries, numerical coincidences, asynchronous confirmations, failed writes,
recoveries, and permitted successor computations are simulated as in that
controller.  The event $W_j=1$ implies an instance of its selected-root
confirmation or marked first-message integrity event.  The existing
uniform-root authentication bound hence gives
\[
 \Pr_{\widetilde V^{[j]}}[W_j=1]
 \le\Omega_j^{\rm auth}.
\]
In particular the shared-compression schedule, the confirmation rows,
fresh-tag guesses, framed-input collisions, and first-message integrity
contribute exactly the event costs displayed in
\eqref{eq:native-auth-continuation}.  Additional eligibility and source
requirements can only restrict this event in the final world.

Each fixed candidate is thus bounded by
$\min\{1,\varepsilon_{{\rm rec},j}+a_j^{\rm auth}+\delta_j
+\Omega_j^{\rm auth}\}$.  Substitute these bounds in
\eqref{eq:native-auth-localization} and cap the result by one.  The partition
was used only in the original native game.  Ideal continuations retain their
own selectors and event flags throughout, so no partition of different
hybrid worlds is required.
\end{proof}

\paragraph{What the composition establishes.}
The theorem connects an eligible native recovery execution to origin
authentication of its selected confirmation schedule and marked first
message.  The comparison uses a single probability event at every stage;
record and primitive reductions therefore enter with coefficient one.  The
operational component advantage $a_j^{\rm auth}$ is attached to this event's
derived distinguisher.  Its value is independently controlled by the same
component security assumption as any output bit admitted by that experiment.
The target scope, final eligibility, and whole-continuation budgets above
specify exactly the executions to which the bound applies.
